% LZSCC.300 Third Year Project -- Project Proposal Template
%
% Compile with pdfLaTeX or LuaLaTeX, and Biber for the references
% (e.g. `latexmk -pdf proposal.tex`; Overleaf does this automatically).
%
% Length: 3-6 pages plus references.
% The proposal is AI RED: no generative AI for any part of the text (see the module handbook).
% The grey guidance text explains what goes into each section. Delete it as you
% write, or hide all of it at once with \guidancefalse below.

\documentclass[11pt,a4paper]{article}

% ----- guidance switch ----------------------------------------------------------
\newif\ifguidance
\guidancetrue % change to \guidancefalse to hide all guidance text

% ----- fonts and language -------------------------------------------------------
\usepackage{iftex}
\ifPDFTeX
	\usepackage[T1]{fontenc}
	\usepackage{lmodern}
\else
	\usepackage{fontspec}
\fi
\usepackage[british]{babel}
\usepackage{csquotes}
\usepackage{microtype}

% ----- layout -------------------------------------------------------------------
\usepackage[a4paper,margin=2.5cm]{geometry}
\usepackage{parskip}
\usepackage{enumitem}

% ----- figures, tables, units -----------------------------------------------------
\usepackage{graphicx}
\usepackage{booktabs}
\usepackage{caption}
\captionsetup[table]{position=top} % table captions go above the table, figure captions below
\usepackage{tabularx}
\usepackage{siunitx}
\usepackage{xcolor}
\usepackage{pgfgantt}

% ----- references -----------------------------------------------------------------
% Use one style consistently: IEEE (numeric, below) or Harvard (author-date: style=authoryear).
\usepackage[style=ieee,backend=biber]{biblatex}
\addbibresource{references.bib}

% ----- links and cross-references (load last) -------------------------------------
\usepackage[hidelinks]{hyperref}
\usepackage[capitalise,noabbrev]{cleveref}

% ----- guidance text --------------------------------------------------------------
\newcommand{\guidance}[1]{\ifguidance{\color{gray}\itshape #1\par}\fi}

% ----- your details -----------------------------------------------------------------
\newcommand{\projecttitle}{Your Project Title}
\newcommand{\studentname}{Your Name}
\newcommand{\studentid}{12345678}
\newcommand{\supervisor}{Your Supervisor}

\hypersetup{pdftitle={\projecttitle},pdfauthor={\studentname}}

\begin{document}

% ----- title block --------------------------------------------------------------------
\begin{center}
	{\small LZSCC.300 Third Year Project -- Project Proposal}\par\medskip
	{\LARGE\bfseries \projecttitle\par}\medskip
	\studentname{} (\studentid) \quad\textbullet\quad Supervisor: \supervisor \quad\textbullet\quad \today
\end{center}

\begin{abstract}
	\guidance{About 150 words, written last. In four to five sentences: the overall aim of the project, why it matters (motivation), how you will do it (methods), and what you expect to find or deliver (expected results).}
\end{abstract}

% ===================================================================================
\section{Introduction}\label{sec:introduction}

\guidance{Motivate the project. Go from general to specific: start with the wider context, narrow down to the problem you address, and back up every claim with evidence and a citation~\cite{example-online}. State the evidence, not who said it: write \enquote{X has increased by 18\%~[3]}, not \enquote{The Guardian reports that \dots}.}

\guidance{End with a short paragraph that summarises your proposed project, and a signposting paragraph: \enquote{The remainder of this proposal is structured as follows. \Cref{sec:background} \dots}}

% ===================================================================================
\section{Background and Related Work}\label{sec:background}

\guidance{Show where your project fits into existing work: similar projects, tools, systems, and papers~\cite{example-article,example-inproceedings}. This is an integrated synthesis, not a list of who did what. For each source, explain why it is relevant to your project, what its limitations are (\enquote{however, \dots}), and how your project addresses them (\enquote{therefore, \dots}).}

\guidance{Organise the section thematically, chronologically, methodologically, or as a funnel (general background $\rightarrow$ closely related work). End with the scope of your project: what exactly you will do, and what not.}

% ===================================================================================
\section{Proposed Project}\label{sec:project}

\subsection{Aims and Objectives}\label{sec:objectives}

\guidance{State the overall aim in one sentence. Then list the objectives. Every objective must be measurable: how will you (and the examiners) know that you achieved it?}

The aim of this project is to \dots

\begin{enumerate}[label=\textbf{O\arabic*}]
	\item \textbf{Objective one} -- \guidance{e.g.\ \enquote{Implement a route planner that finds optimal routes on the map of Leipzig in under one second.}}
	\item \textbf{Objective two} -- \dots
	\item \textbf{Objective three} -- \dots
\end{enumerate}

\subsection{Methodology}\label{sec:methodology}

\guidance{Describe in detail how you will achieve each objective: development process, architecture and technologies, algorithms, datasets, experimental setup, user studies, \dots}

\subsection{Evaluation}\label{sec:evaluation}

\guidance{Your plan for analysis: how will you evaluate the results? Name concrete metrics, baselines, benchmarks, test suites, or study designs, and say how they relate to the objectives in \cref{sec:objectives}.}

\subsection{Ethics}\label{sec:ethics}

\guidance{Only if your project involves people (user studies, surveys, interviews) or personal data: explain what you will collect, how you will protect participants and data, and when you will apply for ethics approval. Otherwise delete this subsection.}

% ===================================================================================
\section{Programme of Work}\label{sec:programme}

\guidance{Break the project into stages and say what each stage delivers and how long it takes. Show the schedule as a Gantt chart (\cref{fig:gantt}). Include time for writing, buffer time, and the module milestones (proposal, assessment forms). Writing can and should start early.}

\begin{figure}[htbp]
	\centering
	\begin{ganttchart}[
			x unit=.55cm,
			y unit title=.5cm,
			title height=1,
			y unit chart=.6cm,
			vgrid={*{9}{gray!25},{gray!80,line width=1pt},*{10}{gray!25}},
			hgrid=false,
			title label font=\footnotesize,
			bar label font=\small,
			milestone label font=\small\itshape,
			bar/.append style={fill=blue!30},
			milestone/.append style={fill=red!70},
		]{1}{20}
		\gantttitle{Michaelmas term (weeks)}{10}
		\gantttitle{Lent term (weeks)}{10} \\
		\gantttitlelist{1,...,20}{1} \\
		\ganttbar{Background research}{1}{4} \\
		\ganttmilestone{Proposal}{3} \\
		\ganttbar{Design}{3}{5} \\
		\ganttbar{Implementation}{5}{13} \\
		\ganttbar{Evaluation}{12}{16} \\
		\ganttbar{Writing the report}{6}{19} \\
		\ganttbar{Buffer}{19}{20} \\
		\ganttmilestone{Assessment forms}{5}
		\ganttmilestone{}{10}
		\ganttmilestone{}{15}
	\end{ganttchart}
	\caption{Project schedule. There is a four-week break between weeks 10 and 11.}
	\label{fig:gantt}
\end{figure}

% ===================================================================================
\section{Resources Required}\label{sec:resources}

\guidance{Hardware, software, licences, datasets, compute, access to people or organisations -- and who provides them. Write \enquote{No special resources are required.} if that is the case.}

% ===================================================================================
\section{Risks and Contingency}\label{sec:risks}

\guidance{What could go wrong, and what will you do then? Think about risks specific to your project.}

\begin{table}[htbp]
	\centering
	\caption{Risks and mitigation.}
	\label{tab:risks}
	\begin{tabularx}{\linewidth}{>{\raggedright\arraybackslash}X l l >{\raggedright\arraybackslash}X}
		\toprule
		Risk & Likelihood & Impact & Mitigation \\
		\midrule
		Dataset not available in time & medium & high & Use synthetic data first, contact provider early \\
		Data loss & low & high & Version control, regular backups \\
		Illness or other absence & low & medium & Buffer time in the schedule \\
		\bottomrule
	\end{tabularx}
\end{table}

% ===================================================================================
\printbibliography[heading=bibintoc]

\end{document}
